% Clase 3 — LPO e Inferencia
% Lectura: AIMA Cap. 8-9
\section{LPO e Inferencia I}

\subsection{Bases de Conocimiento en LPO}

Las bases de conocimiento en LPO pueden expresar relaciones complejas de forma compacta, a diferencia de la lógica proposicional.

\subsection{Unificación}

Para aplicar reglas de inferencia en LPO se necesita \textbf{unificar} expresiones:

$$\text{UNIFY}(\alpha, \beta) = \theta \text{ tal que } \text{SUBST}(\theta, \alpha) = \text{SUBST}(\theta, \beta)$$

Algoritmo: recorre la estructura de ambas expresiones, construyendo sustituciones. Devuelve \textit{fallo} si no hay unificador.

\subsection{Inferencia con Cuantificadores}

\begin{itemize}
    \item \textbf{Instanciación universal}: $\forall x\ \alpha(x) \vdash \alpha(c)$ para cualquier constante $c$.
    \item \textbf{Instanciación existencial}: $\exists x\ \alpha(x) \vdash \alpha(k)$ donde $k$ es una nueva constante de Skolem.
    \item \textbf{Skolemización}: eliminar cuantificadores existenciales introduciendo constantes/funciones de Skolem.
\end{itemize}

\subsection{Forma Normal de Clausura (CNF en LPO)}

Pasos para convertir a CNF:
\begin{enumerate}
    \item Eliminar bicondicionales e implicaciones.
    \item Mover $\neg$ hacia adentro (De Morgan, reglas de cuantificadores).
    \item Estandarizar variables.
    \item Skolemizar.
    \item Eliminar cuantificadores universales.
    \item Distribuir $\lor$ sobre $\land$.
\end{enumerate}
